% Fragmento público generado desde propietarios íntegros.
% Residencia: 52.4.2.
% Fuente: ../incoming/radion_monografia_20260827/manuscrito/sections/05_geometria_accion.tex:218-290
\paragraph{Correspondance entre phase, action, aire et orientation}

\begin{figure}[H]
\centering
\begin{tikzpicture}[scale=1.12]
  \def\aa{4.0}
  \def\bb{2.95}
  \def\cc{2.70}
  \draw[->,RadGray] (-4.7,0)--(4.7,0) node[right]{\(Q\)};
  \draw[->,RadGray] (0,-3.5)--(0,3.5) node[above]{\(P\)};
  \fill[RadTeal!8] (0,0) ellipse[x radius=\aa,y radius=\bb];
  \draw[RadTeal,line width=1.35pt] (0,0) ellipse[x radius=\aa,y radius=\bb];
  \fill[RadCopper] (-\cc,0) circle (2pt) node[below=2mm]{\(F_-\)};
  \fill[RadCopper] (\cc,0) circle (2pt) node[below=2mm]{\(F_+\)};
  \draw[RadGold,line width=1.1pt,-{Latex}] (\aa,0)
     arc[start angle=0,end angle=57.3,x radius=\aa,y radius=\bb];
  \node[RadGold] at (3.7,2.3) {un radion};
  \draw[dashed,RadBlue] (0,0)--(\aa,0) node[midway,below]{\(a_S\)};
  \draw[dashed,RadBlue] (0,0)--(0,\bb) node[midway,left]{\(b_S\)};
  \node[align=center,RadNavy] at (0,-3.25)
    {aire totale \(\pi a_Sb_S=2\pi\hret=\hfull\)};
\end{tikzpicture}
\caption{L'ellipse d'action. Un arc paramétrique d'un radian balaie
\(\hret\) ; un tour enferme \(\hfull\). Les foyers codent l'ouverture
\(C^*\) dans la carte distinguée.}
\label{fig:elipse-accion}
\end{figure}

\paragraph{Relation entre l'action réduite, l'incertitude et l'aire}

La relation \(\hbar=h/(2\pi)\) est souvent présentée comme une commodité
de notation : lorsqu'on décrit des oscillations, des rotations et des transformées de Fourier,
les facteurs \(2\pi\) disparaissent si l'on utilise \(\hbar\). L'ellipse d'action
révèle une lecture plus forte. La division par \(2\pi\) est exactement le
passage de l'action d'un tour à l'action par unité de phase :
\[
h=\mathcal A(2\pi),\qquad \hbar=\mathcal A(1).
\]
La constante réduite n'est pas seulement \(h\) avec une notation plus commode ; c'est la
densité surfacique d'action par rapport à la phase.

Heisenberg, Dirac et Jordan ont établi le langage des variables conjuguées,
des commutateurs et des transformations canoniques. On ne leur attribue pas ici une
\enquote{ellipse ontologique} littérale sans source. La contribution HMT--MD consiste à
identifier la figure positive, à construire ses demi-axes à partir de \((A,C^*)\) et à
démontrer \eqref{eq:area-theorem}. La généalogie historique ouvre le langage ;
le théorème actuel fixe l'objet.

\begin{narrativebox}
L'action ne s'ajoute pas ensuite au mouvement. Dans l'ellipse du radion,
l'action est l'aire que le mouvement génère. Un tour ne transporte pas une
constante préfabriquée : il l'enferme.
\end{narrativebox}

\Needspace{10\baselineskip}
\paragraph{Contrôles algébriques des invariants focaux}

Les ressemblances décimales sont soumises à un contrôle :
\begin{center}
\begin{tabularx}{.96\textwidth}{l r X}
\toprule
Comparaison & Résidu & Verdict\\
\midrule
\(d_1\) face à \(0.54\) & \(4.5455\times10^{-3}\) & pas d'égalité ; le \(54\) ne se déduit pas ;\\
\(d_1^2\) face à \(8/27\) & \(2.3352\times10^{-4}\) & proximité sans application ;\\
\(\eta\) frente a \(3/10\) & \(-3.1032\times10^{-4}\) & no igualdad;\\
\(e_f\) face à \(2/3\) & \(4.7852\times10^{-3}\) & \(2/3\) appartient exactement à \(B_0\) ;\\
\(e^{-\eta}\) face à \(20/27\) & \(3.0740\times10^{-4}\) & no igualdad.\\
\bottomrule
\end{tabularx}
\end{center}
Les invariants solides sont \eqref{eq:eta}, \eqref{eq:excentricidad} et
\eqref{eq:invariante-focal-relativo}. Le critère HMT ne consiste pas à récompenser
une ressemblance décimale, mais à exiger objet, opération et conservation.
